% Editionshinweis zur Übertragung ins heutige Deutsch.
\begin{center}\large Zu dieser Ausgabe\end{center}
\noindent Carl Friedrich Gauss, \emph{Werke}, Band VIII, Göttingen 1900.
Die Königliche Gesellschaft der Wissenschaften zu Göttingen gab die historische Ausgabe heraus.
Zu den historischen Herausgebern und Bearbeitern gehören Robert Fricke, Börsch, Krüger, Stäckel, Weingarten, Brendel und F. Klein.

\noindent Dieser Lesetext enthält die historischen Titelblätter und das Porträt
sowie die vollständige Abteilung Arithmetik und Algebra mit ihren Notizen und
Bemerkungen. Der Lesetext umfasst die gesamte Abteilung Analysis einschließlich der Fragmente über Modulfunktionen und das Pentagramma mirificum sowie die vollständigen Abteilungen Numerisches Rechnen und Wahrscheinlichkeitsrechnung. Die vollständigen Abteilungen Grundlagen der Geometrie, Geometria Situs und Elementare Geometrie umfassen den jeweiligen Nachlass, Briefwechsel, die Anzeigen und die gezeichneten Bemerkungen. Der Lesetext enthält auch die vollständigen Abteilungen über komplexe Größen in der Geometrie und zur Theorie der krummen Flächen sowie die historischen Bemerkungen, Berichtigungen, das Inhaltsverzeichnis und die Verlagsanzeige. Die gezeichneten Bemerkungen behalten ihre historischen Zuschreibungen. Der abgeschnittene rechte Teil der historischen Seite 95 folgt der alternativen Fotografie der Universität Kyūshū. Das Verhältnis der fotografierten Exemplare bleibt ungeklärt. Seitenverweise im Text beziehen sich auf die historische Ausgabe.

\noindent Übertragung ins heutige Deutsch und Satz durch OpenAI GPT-6.1 Sol (Ultra).
Transkription in den Originalsprachen und bearbeitbare mathematische Zeichnungen durch OpenAI GPT-6 Astra (Ultra).
Die unsignierte Titelvignette stammt aus dem \href{https://commons.wikimedia.org/w/index.php?oldid=1179960952}{zweiten Abdruck der \textit{Werke} I von 1870}, digitalisiert nach dem Exemplar der University of California Libraries. Dasselbe Motiv steht auf dem \href{https://gdz.sub.uni-goettingen.de/id/PPN236010751}{Göttinger Titelblatt von 1900}. \href{https://creativecommons.org/publicdomain/mark/1.0/}{Public Domain Mark 1.0}.
\clearpage
